% GENERATED FILE. Edit canonical lessons, then run npm run generate:books.
% canonical-source-hash: fnv1a64:e6b55646c7141a49
% canonical-lessons: LA-C462-nundinor, LA-C462-clamito, LA-C462-vendito, LA-C462-mercimonium, LA-C462-institor, LA-R462-review-455-462

\chapter{To trade at market, To keep shouting, To offer for sale, Goods, Hawker}
\label{ch:la-to-trade-at-market-to-keep-shouting-to-offer-for-sale-goods-hawker}
\hlchaptermodality{462}

\begin{chapteropening}
\textbf{By the end of this chapter:} \emph{I can describe market day: trading at the market, sellers shouting, goods offered for sale, and the hawker.}

Complete the last lesson of chapter 462: I can describe market day: trading at the market, sellers shouting, goods offered for sale, and the hawker.
\end{chapteropening}

\section[nūndinor, nūndinārī]{nūndinor, nūndinārī — to trade at market, to buy and sell}

\label{lesson:LA-C462-nundinor}



\begin{quote}
Before the new one: say the Latin for smoky, then the Latin for filthy.
\end{quote}

\subsection*{You'll want to know: nūndinor, nūndinārī}
\textbf{nūndinor, nūndinārī} — \textquotedblleft{}to trade at market, to buy and sell\textquotedblright{}.

\begin{quote}
Agricolae nōnō quōque diē in oppidō nūndinantur.
\end{quote}

\begin{quote}
\emph{Every ninth day the farmers trade at the market in the town.}
\end{quote}

Deponent: \emph{nūndinor, nūndinārī, nūndinātus sum}. The market day, \emph{nūndinae}, came every ninth day, counting both ends.

\begin{cousinweb}
\emph{nūndinor} is built on \emph{nūndinae} 'market day' (Lewis \& Short), which comes from \emph{novem} 'nine' and \emph{diēs} 'day': the market held every ninth day.
\end{cousinweb}

\begin{grammarlens}[title={What you've built}]
One more word for telling, explaining and describing at B1.
\end{grammarlens}

\subsection*{Your turn}
\begin{itemize}
\raggedright
  \item \emph{Say it:} \emph{nūndinor, nūndinārī}
  \item \emph{Say it:} \emph{nūndinor, nūndinārī}, once more
  \item \emph{Say it:} a sentence of your own with \emph{nūndinor, nūndinārī}
  \item \emph{Recall:} say the Latin for smoky, then the Latin for filthy, then say \emph{nūndinor, nūndinārī} again
\end{itemize}

\begin{tcolorbox}[breakable,colback=teal!4,colframe=teal!35!black,title={Before you move on}]
What does nūndinor, nūndinārī mean? (\textquotedblleft{}To trade at market, to buy and sell\textquotedblright{}.) Say it once more.
\end{tcolorbox}
\section[clāmitō, clāmitāre]{clāmitō, clāmitāre — to keep shouting, to cry out again and again}

\label{lesson:LA-C462-clamito}



\begin{quote}
Before the new one: say the Latin for filthy, then the Latin for to trade at market.
\end{quote}

\subsection*{You'll want to know: clāmitō, clāmitāre}
\textbf{clāmitō, clāmitāre} — \textquotedblleft{}to keep shouting, to cry out again and again\textquotedblright{}.

\begin{quote}
Mercātōrēs in forō clāmitant.
\end{quote}

\begin{quote}
\emph{The merchants in the forum keep shouting.}
\end{quote}

\emph{clāmitō, clāmitāre, clāmitāvī, clāmitātum}, a frequentative of \emph{clāmō}: the \emph{-itō} means 'again and again'.

\begin{cousinweb}
\emph{clāmitō} is built on \emph{clāmō} 'to shout' with the frequentative ending \emph{-itō} (EtymDB; Lewis \& Short).
\end{cousinweb}

\begin{grammarlens}[title={What you've built}]
One more word for telling, explaining and describing at B1.
\end{grammarlens}

\subsection*{Your turn}
\begin{itemize}
\raggedright
  \item \emph{Say it:} \emph{clāmitō, clāmitāre}
  \item \emph{Say it:} \emph{clāmitō, clāmitāre}, once more
  \item \emph{Say it:} a sentence of your own with \emph{clāmitō, clāmitāre}
  \item \emph{Recall:} say the Latin for filthy, then the Latin for to trade at market, then say \emph{clāmitō, clāmitāre} again
\end{itemize}

\begin{tcolorbox}[breakable,colback=teal!4,colframe=teal!35!black,title={Before you move on}]
What does clāmitō, clāmitāre mean? (\textquotedblleft{}To keep shouting, to cry out again and again\textquotedblright{}.) Say it once more.
\end{tcolorbox}
\section[vēnditō, vēnditāre]{vēnditō, vēnditāre — to offer for sale, to try to sell}

\label{lesson:LA-C462-vendito}



\begin{quote}
Before the new one: say the Latin for to trade at market, then the Latin for to keep shouting.
\end{quote}

\subsection*{You'll want to know: vēnditō, vēnditāre}
\textbf{vēnditō, vēnditāre} — \textquotedblleft{}to offer for sale, to try to sell\textquotedblright{}.

\begin{quote}
Mercātor vīnum suum omnibus vēnditat.
\end{quote}

\begin{quote}
\emph{The merchant tries to sell his wine to everyone.}
\end{quote}

\emph{vēnditō, vēnditāre, vēnditāvī, vēnditātum}, a frequentative of \emph{vendō}: offering again and again, praising your goods.

\begin{cousinweb}
\emph{vēnditō} is built on \emph{vendō} 'to sell' with the frequentative ending \emph{-itō} (Lewis \& Short). English \emph{vend} and \emph{vendor} come from \emph{vendō}.
\end{cousinweb}

\begin{grammarlens}[title={What you've built}]
One more word for telling, explaining and describing at B1.
\end{grammarlens}

\subsection*{Your turn}
\begin{itemize}
\raggedright
  \item \emph{Say it:} \emph{vēnditō, vēnditāre}
  \item \emph{Say it:} \emph{vēnditō, vēnditāre}, once more
  \item \emph{Say it:} a sentence of your own with \emph{vēnditō, vēnditāre}
  \item \emph{Recall:} say the Latin for to trade at market, then the Latin for to keep shouting, then say \emph{vēnditō, vēnditāre} again
\end{itemize}

\begin{tcolorbox}[breakable,colback=teal!4,colframe=teal!35!black,title={Before you move on}]
What does vēnditō, vēnditāre mean? (\textquotedblleft{}To offer for sale, to try to sell\textquotedblright{}.) Say it once more.
\end{tcolorbox}
\section[mercimōnium, mercimōniī]{mercimōnium, mercimōniī — goods, wares, merchandise}

\label{lesson:LA-C462-mercimonium}



\begin{quote}
Before the new one: say the Latin for to keep shouting, then the Latin for to offer for sale.
\end{quote}

\subsection*{You'll want to know: mercimōnium, mercimōniī}
\textbf{mercimōnium, mercimōniī} — \textquotedblleft{}goods, wares, merchandise\textquotedblright{}.

\begin{quote}
Mercātor mercimōnia in nāve ad urbem vehit.
\end{quote}

\begin{quote}
\emph{The merchant carries his goods by ship to the city.}
\end{quote}

Second declension, neuter: \emph{mercimōnium, mercimōniī n.}, mostly plural. A word of comedy, from \emph{merx} 'goods'.

\begin{cousinweb}
\emph{mercimōnium} is built on \emph{merx, mercis} 'merchandise' with the suffix \emph{-mōnium} (Lewis \& Short). English \emph{merchant} and \emph{commerce} come from the same family.
\end{cousinweb}

\begin{grammarlens}[title={What you've built}]
One more word for telling, explaining and describing at B1.
\end{grammarlens}

\subsection*{Your turn}
\begin{itemize}
\raggedright
  \item \emph{Say it:} \emph{mercimōnium, mercimōniī}
  \item \emph{Say it:} \emph{mercimōnium, mercimōniī}, once more
  \item \emph{Say it:} a sentence of your own with \emph{mercimōnium, mercimōniī}
  \item \emph{Recall:} say the Latin for to keep shouting, then the Latin for to offer for sale, then say \emph{mercimōnium, mercimōniī} again
\end{itemize}

\begin{tcolorbox}[breakable,colback=teal!4,colframe=teal!35!black,title={Before you move on}]
What does mercimōnium, mercimōniī mean? (\textquotedblleft{}Goods, wares, merchandise\textquotedblright{}.) Say it once more.
\end{tcolorbox}
\section[īnstitor, īnstitōris]{īnstitor, īnstitōris — a hawker, a pedlar; a shop agent}

\label{lesson:LA-C462-institor}



\begin{quote}
Before the new one: say the Latin for to offer for sale, then the Latin for goods.
\end{quote}

\subsection*{You'll want to know: īnstitor, īnstitōris}
\textbf{īnstitor, īnstitōris} — \textquotedblleft{}a hawker, a pedlar; a shop agent\textquotedblright{}.

\begin{quote}
Īnstitor per viās it et mercēs suās clāmat.
\end{quote}

\begin{quote}
\emph{The hawker goes through the streets and calls out his goods.}
\end{quote}

Third declension, masculine: \emph{īnstitor, īnstitōris m.} A man who sells for someone else, often on the street.

\begin{cousinweb}
\emph{īnstitor} is built on \emph{īnsistō} 'to stand on, to keep at a thing' (Lewis \& Short), from \emph{in-} and \emph{sistō} 'to stand'. EtymDB lists it without a parent.
\end{cousinweb}

\begin{grammarlens}[title={What you've built}]
One more word for telling, explaining and describing at B1.
\end{grammarlens}

\subsection*{Your turn}
\begin{itemize}
\raggedright
  \item \emph{Say it:} \emph{īnstitor, īnstitōris}
  \item \emph{Say it:} \emph{īnstitor, īnstitōris}, once more
  \item \emph{Say it:} a sentence of your own with \emph{īnstitor, īnstitōris}
  \item \emph{Recall:} say the Latin for to offer for sale, then the Latin for goods, then say \emph{īnstitor, īnstitōris} again
\end{itemize}

\begin{tcolorbox}[breakable,colback=teal!4,colframe=teal!35!black,title={Before you move on}]
What does īnstitor, īnstitōris mean? (\textquotedblleft{}A hawker, a pedlar; a shop agent\textquotedblright{}.) Say it once more.
\end{tcolorbox}
\section[(dialogue)]{Words from chapters 455–462 — a review}

\label{lesson:LA-R462-review-455-462}



\begin{quote}
Say the Latin for to catch sight of and to come near.
\end{quote}

\subsection*{Your turn}
Say these aloud:

\begin{itemize}
\raggedright
  \item to catch sight of, to come near, to turn back, to delay, to stop for a while
  \item a tenant, a townsman, a resident, to live in, an occupant
  \item a block of flats, an upper room, a storey, a timber, a beam
  \item an alley, a neighbourhood, a latrine, a sewer, a water tower
  \item a colonnade, a temple, a shrine, the sacred boundary, to build up
  \item a crowd, crowded, to throng, to go back and forth, near the city
  \item a hot-drinks bar, hot water, to drink, smoky, filthy
  \item to trade at market, to keep shouting, to offer for sale, goods, a hawker
\end{itemize}

\begin{tcolorbox}[breakable,colback=teal!4,colframe=teal!35!black,title={Before you move on}]
To catch sight of? (\textbf{cōnspicor, cōnspicārī}.) To come near? (\textbf{accēdō, accēdere}.) To turn back? (\textbf{revertor, revertī}.) To delay? (\textbf{moror, morārī}.) To stop for a while? (\textbf{commoror, commorārī}.) A tenant? (\textbf{inquilīnus, inquilīnī}.) A townsman? (\textbf{oppidānus, oppidānī}.) A resident? (\textbf{incola, incolae}.) To live in? (\textbf{incolō, incolere}.) An occupant? (\textbf{habitātor, habitātōris}.) A block of flats? (\textbf{īnsula, īnsulae}.) An upper room? (\textbf{cēnāculum, cēnāculī}.) A storey? (\textbf{contignātiō, contignātiōnis}.) A timber? (\textbf{tignum, tignī}.) A beam? (\textbf{trabs, trabis}.) An alley? (\textbf{angiportum, angiportī}.) A neighbourhood? (\textbf{vīcīnia, vīcīniae}.) A latrine? (\textbf{lātrīna, lātrīnae}.) A sewer? (\textbf{cloāca, cloācae}.) A water tower? (\textbf{castellum, castellī}.) A colonnade? (\textbf{porticus, porticūs}.) A temple? (\textbf{aedēs, aedis}.) A shrine? (\textbf{aedicula, aediculae}.) The sacred boundary? (\textbf{pōmērium, pōmēriī}.) To build up? (\textbf{exstruō, exstruere}.) A crowd? (\textbf{turba, turbae}.) Crowded? (\textbf{frequēns, frequentis}.) To throng? (\textbf{frequentō, frequentāre}.) To go back and forth? (\textbf{commeō, commeāre}.) Near the city? (\textbf{suburbānus, suburbāna, suburbānum}.) A hot-drinks bar? (\textbf{thermopōlium, thermopōliī}.) Hot water? (\textbf{calda, caldae}.) To drink? (\textbf{pōtō, pōtāre}.) Smoky? (\textbf{fūmōsus, fūmōsa, fūmōsum}.) Filthy? (\textbf{immundus, immunda, immundum}.) To trade at market? (\textbf{nūndinor, nūndinārī}.) To keep shouting? (\textbf{clāmitō, clāmitāre}.) To offer for sale? (\textbf{vēnditō, vēnditāre}.) Goods? (\textbf{mercimōnium, mercimōniī}.) A hawker? (\textbf{īnstitor, īnstitōris}.)
\end{tcolorbox}
